% segment-local-id: OR03-S010
% segment-id: d90.orig.v1.tr03.seg0010
% license: CC BY-SA 4.0
% provenance: Materi asesmen asli berbahasa Indonesia, disusun secara independen.

\section{Ujian tengah kumulatif}
\label{orig03:midterm}

% stable-id: d90.orig.v1.tr03.midterm.exercise.0001
\begin{exercise}[Keberadaan dan ketunggalan model kuadrat--$\ell^1$]
\label{orig03:midterm:exercise:0001}
Untuk $A\in\mathbb R^{m\times n}$, $b\in\mathbb R^m$, dan $\lambda>0$,
pertimbangkan
\[
 F(x)=\lVert Ax-b\rVert_2^2+\lambda\lVert x\rVert_1.
\]
Buktikan bahwa $F$ selalu mempunyai peminimum. Apakah peminimum itu selalu
tunggal? Berikan satu kondisi sederhana pada $A$ yang menjamin ketunggalan.
\end{exercise}

% stable-id: d90.orig.v1.tr03.midterm.hint.0001
\paragraph{Petunjuk bertahap.}
\label{orig03:midterm:hint:0001}
\textbf{Tahap 1.} Bandingkan $F(x)$ dengan $\lambda\lVert x\rVert_1$ ketika
$\lVert x\rVert_2\to\infty$.
\textbf{Tahap 2.} Untuk ketunggalan, cari kondisi yang membuat suku kuadrat
konveks ketat; untuk menyangkal ketunggalan umum, izinkan $A$ memiliki
kernel.

% stable-id: d90.orig.v1.tr03.midterm.answer.0001
\paragraph{Jawaban singkat.}
\label{orig03:midterm:answer:0001}
$F$ kontinu dan koersif karena
$F(x)\ge\lambda\lVert x\rVert_1$, sehingga minimum dicapai. Ketunggalan tidak
selalu berlaku; misalnya $A=0$ dan $b=0$ sebenarnya memberi minimum tunggal,
tetapi matriks dengan kolom identik dapat menghasilkan banyak representasi
$\ell^1$ dengan nilai sama. Kondisi sederhana yang cukup adalah $A$ berperingkat
kolom penuh.

% stable-id: d90.orig.v1.tr03.midterm.solution.0001
\paragraph{Solusi lengkap.}
\label{orig03:midterm:solution:0001}
Kedua suku $F$ kontinu dan tak negatif. Karena
$\lVert x\rVert_1\ge\lVert x\rVert_2$,
\[
 F(x)\ge\lambda\lVert x\rVert_2\longrightarrow\infty
 \quad\text{ketika }\lVert x\rVert_2\to\infty.
\]
Jadi semua himpunan sublevel tak kosong yang relevan kompak, dan teorema
Weierstrass memberi keberadaan peminimum.

Ketunggalan dapat gagal. Sebagai contoh, ambil
$A=(1\;1)$, $b=1$, dan $0<\lambda<2$. Untuk $x_1,x_2\ge0$ dengan
$x_1+x_2=s$, objektif hanya bergantung pada $s$:
$(s-1)^2+\lambda s$. Nilai minimumnya dicapai pada
$s=1-\lambda/2>0$, dan setiap pembagian nonnegatif dari $s$ memberi
peminimum berbeda. Jika $A$ berperingkat kolom penuh, maka
$x\mapsto\lVert Ax-b\rVert^2$ konveks ketat. Jumlah fungsi konveks ketat
dengan fungsi konveks tetap konveks ketat, sehingga peminimum yang sudah
terbukti ada harus tunggal.

% stable-id: d90.orig.v1.tr03.midterm.exercise.0002
\begin{exercise}[Konjugat kuadrat bergeser]
\label{orig03:midterm:exercise:0002}
Untuk $\mu>0$ dan $c\in\mathbb R^n$, definisikan
\[
 f(x)=\frac\mu2\lVert x\rVert^2+\langle c,x\rangle.
\]
Hitung $f^*$ secara langsung dari definisi, verifikasi $f^{**}=f$, dan
tentukan konstanta Lipschitz terkecil dari $\nabla f^*$.
\end{exercise}

% stable-id: d90.orig.v1.tr03.midterm.hint.0002
\paragraph{Petunjuk bertahap.}
\label{orig03:midterm:hint:0002}
\textbf{Tahap 1.} Maksimalkan kuadrat konkaf dalam definisi
$f^*(y)=\sup_x\{\langle y,x\rangle-f(x)\}$.
\textbf{Tahap 2.} Terapkan perhitungan yang sama pada $f^*$ atau substitusikan
gradien inversnya.

% stable-id: d90.orig.v1.tr03.midterm.answer.0002
\paragraph{Jawaban singkat.}
\label{orig03:midterm:answer:0002}
\[
 f^*(y)=\frac1{2\mu}\lVert y-c\rVert^2,
 \qquad
 \nabla f^*(y)=\frac{y-c}{\mu}.
\]
Karena itu $f^{**}=f$, dan konstanta Lipschitz terkecil gradien $f^*$ adalah
$1/\mu$.

% stable-id: d90.orig.v1.tr03.midterm.solution.0002
\paragraph{Solusi lengkap.}
\label{orig03:midterm:solution:0002}
Definisi konjugat memberi
\[
 f^*(y)=\sup_x\left\{\langle y-c,x\rangle-\frac\mu2\lVert x\rVert^2\right\}.
\]
Titik stasionernya $x=(y-c)/\mu$, dan substitusi menghasilkan
$\lVert y-c\rVert^2/(2\mu)$. Selanjutnya,
\[
 f^{**}(x)=\sup_y\left\{\langle x,y\rangle
 -\frac1{2\mu}\lVert y-c\rVert^2\right\}.
\]
Peminimum negatif kuadrat, atau kondisi stasioner, memberi
$y=c+\mu x$. Nilainya
$\langle c,x\rangle+\mu\lVert x\rVert^2/2=f(x)$. Akhirnya,
\[
 \lVert\nabla f^*(y)-\nabla f^*(z)\rVert
 =\frac1\mu\lVert y-z\rVert.
\]
Kesetaraan berlaku untuk setiap pasangan berbeda, jadi $1/\mu$ bukan hanya
batas yang cukup, melainkan konstanta terkecil.

% stable-id: d90.orig.v1.tr03.midterm.exercise.0003
\begin{exercise}[Ketaksamaan satu langkah gradien proksimal]
\label{orig03:midterm:exercise:0003}
Misalkan $f$ konveks dengan gradien $L$-Lipschitz dan $r$ konveks proper
semikontinu-bawah. Untuk $0<\eta\le1/L$, definisikan
\[
 x^+=\operatorname{prox}_{\eta r}(x-\eta\nabla f(x)),
 \qquad F=f+r.
\]
Buktikan bahwa untuk setiap $u\in\operatorname{dom}r$,
\[
 F(x^+)-F(u)
 \le\frac{\lVert x-u\rVert^2-\lVert x^+-u\rVert^2}{2\eta}.
\]
Kemudian turunkan batas $O(1/k)$ untuk nilai fungsi iterasi gradien proksimal
bila peminimum $x^\star$ ada.
\end{exercise}

% stable-id: d90.orig.v1.tr03.midterm.hint.0003
\paragraph{Petunjuk bertahap.}
\label{orig03:midterm:hint:0003}
\textbf{Tahap 1.} Gabungkan lemma penurunan untuk $f$ dengan inklusi
optimalitas proksimal dan ketaksamaan subgradien untuk $r$.
\textbf{Tahap 2.} Gunakan identitas tiga norma, tetapkan $u=x^\star$, lalu
jumlahkan dan manfaatkan monotoni nilai objektif.

% stable-id: d90.orig.v1.tr03.midterm.answer.0003
\paragraph{Jawaban singkat.}
\label{orig03:midterm:answer:0003}
Suku tambahan yang muncul adalah
$-(1/(2\eta)-L/2)\lVert x^+-x\rVert^2\le0$, sehingga ketaksamaan yang diminta
berlaku. Untuk iterasi $x_{j+1}$ yang sama,
\[
 F(x_k)-F(x^\star)
 \le\frac{\lVert x_0-x^\star\rVert^2}{2\eta k}.
\]

% stable-id: d90.orig.v1.tr03.midterm.solution.0003
\paragraph{Solusi lengkap.}
\label{orig03:midterm:solution:0003}
Inklusi proksimal menyediakan $s^+\in\partial r(x^+)$ dengan
\[
 s^+=-\nabla f(x)-\frac1\eta(x^+-x).
\]
Lemma penurunan, konveksitas $f$, dan ketaksamaan subgradien memberi
\begin{align*}
 F(x^+)-F(u)
 &\le \langle\nabla f(x),x^+-u\rangle
      +\frac L2\lVert x^+-x\rVert^2
      +\langle s^+,x^+-u\rangle\\
 &=-\frac1\eta\langle x^+-x,x^+-u\rangle
      +\frac L2\lVert x^+-x\rVert^2\\
 &=\frac{\lVert x-u\rVert^2-\lVert x^+-u\rVert^2}{2\eta}
   -\left(\frac1{2\eta}-\frac L2\right)\lVert x^+-x\rVert^2.
\end{align*}
Koefisien terakhir nonnegatif karena $\eta\le1/L$, sehingga suku negatif itu
dapat dibuang. Dengan $u=x^\star$, penjumlahan untuk $j=0,\ldots,k-1$
meneleskopkan jarak dan memberi
\[
 \sum_{j=0}^{k-1}\bigl(F(x_{j+1})-F(x^\star)\bigr)
 \le\frac{\lVert x_0-x^\star\rVert^2}{2\eta}.
\]
Ketaksamaan yang sama dengan $u=x_j$ menunjukkan nilai $F(x_j)$ tidak naik.
Karena itu suku terakhir tidak melebihi rata-rata $k$ suku, yang menghasilkan
batas pada jawaban singkat.

% stable-id: d90.orig.v1.tr03.midterm.exercise.0004
\begin{exercise}[Newton dan penerimaan langkah penuh Armijo]
\label{orig03:midterm:exercise:0004}
Misalkan
\[
 f(x)=\tfrac12x^\top Qx-b^\top x,
 \qquad Q\succ0.
\]
Dari titik sembarang $x$, tentukan arah Newton $p$, buktikan bahwa $x+p$
adalah peminimum eksak, dan tunjukkan bahwa uji Armijo
\[
 f(x+p)\le f(x)+c\nabla f(x)^\top p
\]
menerima langkah penuh untuk setiap $0<c\le1/2$.
\end{exercise}

% stable-id: d90.orig.v1.tr03.midterm.hint.0004
\paragraph{Petunjuk bertahap.}
\label{orig03:midterm:hint:0004}
\textbf{Tahap 1.} Selesaikan $Qp=-\nabla f(x)$ dan hitung gradien di $x+p$.
\textbf{Tahap 2.} Gunakan ekspansi kuadrat eksak sepanjang $p$ serta
$\nabla f(x)^\top p=-p^\top Qp$.

% stable-id: d90.orig.v1.tr03.midterm.answer.0004
\paragraph{Jawaban singkat.}
\label{orig03:midterm:answer:0004}
$p=Q^{-1}(b-Qx)$ dan $x+p=Q^{-1}b=x^\star$. Selain itu,
\[
 f(x+p)=f(x)-\tfrac12p^\top Qp,
 \qquad \nabla f(x)^\top p=-p^\top Qp,
\]
sehingga Armijo berlaku tepat bila $c\le1/2$ (atau trivial bila $p=0$).

% stable-id: d90.orig.v1.tr03.midterm.solution.0004
\paragraph{Solusi lengkap.}
\label{orig03:midterm:solution:0004}
Gradiennya $Qx-b$ dan Hessiannya $Q$. Persamaan Newton
$Qp=-(Qx-b)$ memberi $p=Q^{-1}b-x$, sehingga titik baru adalah
$Q^{-1}b$. Karena $Q$ definit positif, titik stasioner ini peminimum tunggal.
Ekspansi eksak memberi
\[
 f(x+p)=f(x)+\nabla f(x)^\top p+\tfrac12p^\top Qp
 =f(x)-\tfrac12p^\top Qp.
\]
Ruas kanan uji Armijo adalah
$f(x)-c\,p^\top Qp$. Untuk $p\ne0$, definit positifnya $Q$ membuat uji
ekuivalen dengan $1/2\ge c$. Jika $p=0$, titik sudah optimal dan kedua ruas
sama untuk semua $c$.

% stable-id: d90.orig.v1.tr03.midterm.exercise.0005
\begin{exercise}[KKT untuk proyeksi pada bola]
\label{orig03:midterm:exercise:0005}
Ambil $R>0$ dan $c\in\mathbb R^n$ dengan $\lVert c\rVert>R$. Selesaikan
\[
 \min_x\ \tfrac12\lVert x-c\rVert^2
 \quad\text{dengan syarat}\quad
 h(x)=\tfrac12(\lVert x\rVert^2-R^2)\le0.
\]
Verifikasi kondisi Slater, tulis seluruh kondisi KKT, dan hitung pasangan
primal--dual optimal secara eksplisit.
\end{exercise}

% stable-id: d90.orig.v1.tr03.midterm.hint.0005
\paragraph{Petunjuk bertahap.}
\label{orig03:midterm:hint:0005}
\textbf{Tahap 1.} Titik nol adalah titik Slater; karena $c$ di luar bola,
kendala pada solusi harus aktif.
\textbf{Tahap 2.} Stasioneritas memberi $(1+\lambda)x=c$.

% stable-id: d90.orig.v1.tr03.midterm.answer.0005
\paragraph{Jawaban singkat.}
\label{orig03:midterm:answer:0005}
KKT terdiri dari
\[
 x-c+\lambda x=0,quad h(x)\le0,quad\lambda\ge0,quad\lambda h(x)=0.
\]
Solusinya
\[
 x^\star=R\frac{c}{\lVert c\rVert},
 \qquad \lambda^\star=\frac{\lVert c\rVert}{R}-1.
\]

% stable-id: d90.orig.v1.tr03.midterm.solution.0005
\paragraph{Solusi lengkap.}
\label{orig03:midterm:solution:0005}
Karena $h(0)=-R^2/2<0$, Slater berlaku. Masalah konveks dan objektifnya
konveks ketat, sehingga KKT perlu dan cukup serta solusi primal tunggal.
Lagrangiannya
\[
 L(x,\lambda)=\tfrac12\lVert x-c\rVert^2
 +\frac\lambda2(\lVert x\rVert^2-R^2).
\]
Stasioneritas dan tiga kondisi lainnya tercantum pada jawaban singkat. Solusi
tanpa kendala $x=c$ tidak layak, maka kendala aktif dan
$\lVert x\rVert=R$. Dari $(1+\lambda)x=c$ diperoleh bahwa $x$ searah dengan
$c$. Menormakan persamaan itu menghasilkan
$1+\lambda=\lVert c\rVert/R$, sehingga pasangan yang dinyatakan mengikuti.
Pengalinya positif, maka semua kondisi KKT terpenuhi.

% stable-id: d90.orig.v1.tr03.midterm.exercise.0006
\begin{exercise}[Ketika kualifikasi gagal dan pengali tidak ada]
\label{orig03:midterm:exercise:0006}
Pertimbangkan masalah konveks skalar
\[
 \min_{x\in\mathbb R} x
 \quad\text{dengan syarat}\quad x^2\le0.
\]
Tentukan solusi, tunjukkan bahwa Slater gagal, dan buktikan bahwa tidak ada
pengali $\lambda\ge0$ yang memenuhi kondisi KKT pada solusi. Jelaskan mengapa
contoh ini tidak menyangkal teorema KKT dengan kualifikasi.
\end{exercise}

% stable-id: d90.orig.v1.tr03.midterm.hint.0006
\paragraph{Petunjuk bertahap.}
\label{orig03:midterm:hint:0006}
\textbf{Tahap 1.} Tentukan seluruh himpunan layak dari ketaksamaan kuadrat.
\textbf{Tahap 2.} Diferensiasikan $L(x,\lambda)=x+\lambda x^2$ di titik layak.

% stable-id: d90.orig.v1.tr03.midterm.answer.0006
\paragraph{Jawaban singkat.}
\label{orig03:midterm:answer:0006}
Satu-satunya titik layak dan optimal ialah $x^\star=0$. Tidak ada titik dengan
$x^2<0$, jadi Slater gagal. Stasioneritas KKT menuntut
$1+2\lambda x^\star=0$, yakni $1=0$, mustahil untuk setiap $\lambda$.

% stable-id: d90.orig.v1.tr03.midterm.solution.0006
\paragraph{Solusi lengkap.}
\label{orig03:midterm:solution:0006}
Karena kuadrat selalu nonnegatif, syarat $x^2\le0$ ekuivalen dengan $x=0$.
Maka $x^\star=0$ adalah satu-satunya titik layak dan otomatis peminimum.
Slater memerlukan titik dengan $x^2<0$, yang tidak ada. Untuk pengali
$\lambda\ge0$, Lagrangian adalah $L(x,\lambda)=x+\lambda x^2$ dan kondisi
stasioneritas pada $x^\star$ ialah
\[
 0=\frac{\partial L}{\partial x}(0,\lambda)=1+2\lambda\cdot0=1,
\]
sebuah kontradiksi. Teorema KKT yang menjamin keberadaan pengali memakai
kondisi kualifikasi seperti Slater. Hipotesis itu gagal di sini, sehingga
ketiadaan pengali justru menunjukkan mengapa kualifikasi tidak boleh
dihilangkan.
